\documentclass[UTF8]{ctexart}
\usepackage{xeCJK}
\usepackage{geometry}
\geometry{left=2cm,right=2cm,top=2.5cm,bottom=2.5cm}
\setlength{\headheight}{12.7pt}

\usepackage{titlesec}
\titleformat{\section}
  {\normalfont\large\bfseries}
  {\thesection}
  {1em}
  {}
\titlespacing*{\section}
  {0pt}
  {3.5ex plus 1ex minus .2ex}
  {2.3ex plus .2ex}

\usepackage{amsmath}
\usepackage{amssymb}
\usepackage{amsthm}
\usepackage{enumitem}
\setlist[enumerate]{itemsep=0pt,topsep=0pt,parsep=0pt,leftmargin=0pt,itemindent=2.5em}
\setlist[itemize]{itemsep=0pt,topsep=0pt,parsep=0pt,leftmargin=0pt,itemindent=2.5em}
\usepackage{fancyhdr}
\pagestyle{fancy}
\fancyhf{}
\usepackage{graphicx}
\usepackage{hyperref}
\usepackage{indentfirst}
\usepackage{lmodern}


\begin{document}
\setlength{\abovedisplayskip}{1pt}
\setlength{\belowdisplayskip}{1pt}

\section{滤子}

% 要学点集拓扑, 不得不先把滤子和网的一些概念搞清楚.

% 周杰伦新专辑还有8分钟就要正式发行, 非常 ~ 非常期待!
% 2026.03.24

\hypertarget{sec:定义1.1}{}
\noindent\textbf{定义1.1 (滤子)}

给定集合$X$和它的子集族$F\subset\mathcal{P}(X)$. 如果:
\begin{enumerate}
    \item $\varnothing\notin F,\,X\in F$,
    \item 对任意的$a_1,a_2\subset X$, 如果$a_1\in F$且$a_1\subset a_2$,
    那么$a_2\in F$,
    \item 对任意的$a_1,a_2\in F$, 有$a_1\cap a_2\in F$,
\end{enumerate}
就称$F$是$X$上的一个\textbf{滤子}.

% 爽听! 今天的工作就到这里.

% 好听!

% 2026.03.25

\vspace{10pt}

\hypertarget{sec:例1.1}{}
\noindent\textbf{例1.1}

给定集合$X$和非空的$b\subset X$, 包含$b$的全体子集
\[
    F_b=\{a\subset X\mid a\supset b\}
\]
是$X$上的滤子. 特别地,
\begin{enumerate}
    \item $F_X=\{X\}$称为$X$上的\textbf{平凡滤子},
    \item 对任意的$x\in X$, $F_{\{x\}}=\{a\subset X\mid x\in a\}$
    称为$x$的\textbf{主滤子}.
\end{enumerate}

\noindent{\kaishu 证明.}

\begin{enumerate}
    \item $\varnothing\not\supset b,\,X\supset b$,
    \item 对任意的$a_1,a_2\subset X$, 如果$a_1\in F_b$且$a_1\subset a_2$,
    那么$a_2\supset a_1\supset b$, 所以$a_2\in F_b$,
    \item 对任意的$a_1,a_2\in F_b$, 有$a_1,a_2\supset b$,
    所以$a_1\cap a_2\supset b$. \hfill $\qedsymbol$
\end{enumerate}

\vspace{10pt}

\hypertarget{sec:例1.2}{}
\noindent\textbf{例1.2 (Fréchet滤子)}

给定无限集合$X$, 称它的余有限子集的全体为$X$上的\textbf{Fréchet滤子}:
\[
    F_\sigma=\{a\subset X\mid a^c\text{\,有限}\}.
\]

\noindent{\kaishu 证明.}

\begin{enumerate}
    \item $\varnothing^c=X$是无限集, $X^c=\varnothing$是有限集,
    \item 对任意的$a_1,a_2\subset X$, 如果$a_1\in F_\sigma$且$a_1\subset a_2$,
    那么$a_2^c\subset a_1^c$也有限,
    \item 对任意的$a_1,a_2\in F_\sigma$, 有$a_1^c,a_2^c$有限, 所以
    $(a_1\cap a_2)^c=a_1^c\cup a_2^c$也有限. \hfill $\qedsymbol$
\end{enumerate}

\vspace{10pt}

\hypertarget{sec:定理1.1}{}
\noindent\textbf{定理1.1 (滤子对有限交封闭)}

给定$X$上的滤子$F$. 对$F$中任意有限个$(a_i)_{i<n}$,
有$\displaystyle\bigcap_{i=0}^{n-1}a_i\in F$.

\noindent{\kaishu 证明.}

首先, $n=1$的情况显然; 其次, 假设任取$(a_i)_{i<n}$
有$\displaystyle\bigcap_{i=0}^{n-1}a_i\in F$, 那么任取$(a_i)_{i<n+1}$有:
\vspace{3pt}
\begin{equation*}
    \bigcap_{i=0}^na_i=\left(\bigcap_{i=0}^{n-1}a_i\right)\cap a_n\in F.
\tag*{\qedsymbol}\end{equation*}





\newpage

\hypertarget{sec:定理1.2}{}
\noindent\textbf{定理1.2}

给定集合$X$和$X$上的非空的一族滤子$\mathcal{S}$ :
\vspace{3pt}
\begin{enumerate}
    \item $\displaystyle\bigcap\mathcal{S}$也是$X$上的滤子,
    \vspace{3pt}
    \item 如果$(\mathcal{S},\subset)$是有向集,
    那么$\displaystyle\bigcup\mathcal{S}$也是$X$上的滤子.
\end{enumerate}

\noindent{\kaishu 证明.}

\begin{enumerate}
    \item 显然$\displaystyle\varnothing\notin\bigcap\mathcal{S},
    \,X\in\bigcap\mathcal{S}$.

    \vspace{3pt}\hspace{2.17em}
    如果$\displaystyle a_1\in\bigcap\mathcal{S}$且$a_1\subset a_2\subset X$, 那么
    对任意的$F\in\mathcal{S}$有
    \[
        a_1\in F\quad\implies\quad a_2\in F,
    \]
    从而$\displaystyle a_2\in\bigcap\mathcal{S}$.

    \vspace{3pt}\hspace{2.17em}
    如果$\displaystyle a_1,a_2\in\bigcap\mathcal{S}$, 那么
    对任意的$F\in\mathcal{S}$有
    \[
        a_1,a_2\in F\quad\implies\quad a_1\cap a_2\in F,
    \]
    从而$\displaystyle a_1\cap a_2\in\bigcap\mathcal{S}$.

    \vspace{3pt}
    \item 显然$\displaystyle\varnothing\notin\bigcup\mathcal{S},
    \,X\in\bigcup\mathcal{S}$.

    \vspace{3pt}\hspace{2.17em}
    如果$\displaystyle a_1\in\bigcup\mathcal{S}$且$a_1\subset a_2\subset X$, 那么
    存在$F\in\mathcal{S}$使得
    \[
        a_1\in F\quad\implies\quad a_2\in F,
    \]
    从而$\displaystyle a_2\in\bigcup\mathcal{S}$.
    
    \vspace{3pt}\hspace{2.17em}
    如果$\displaystyle a_1,a_2\in\bigcup\mathcal{S}$, 那么
    存在$F_1,F_2\in\mathcal{S}$使得
    \[
        a_1\in F_1\quad\land\quad a_2\in F_2,
    \]
    如果$(\mathcal{S},\subset)$是有向集, 那么可以取$F\in\mathcal{S}$
    使得$F_1,F_2\subset F$, 此时
    \[
        a_1,a_2\in F\quad\implies\quad a_1\cap a_2\in F,
    \]
    从而$\displaystyle a_1\cap a_2\in\bigcup\mathcal{S}$. \hfill $\qedsymbol$
\end{enumerate}

\end{document}